\section{Números ordinais}
Vimos que cada natural $0, 1, 2, \dots$ é um conjunto bem ordenado.
Se $n$ é um número natural, $S(n) = n\cup\{n\}$ é também um número natural, e portanto um conjunto bem ordenado.
O conjunto $\omega$ é um conjunto bem ordenado.
A ordem de todos esses conjuntos é dada pela relação de pertencimento $\in$.

Vimos também que toda boa ordem finita é isomorfa a um único número natural.
Assim, os números naturais servem como representantes de todas as boas ordens finitas.

Gostaríamos de encontrar representantes de todas as boas ordens, independente de serem infinitas ou não.
A tentativa natural é seguir executando o mesmo processo de construção dos números naturais, ou seja, continuar a sequência:

\[
0, 1, 2, \dots, \omega, S(\omega), S(S(\omega)), \dots
\]

Podemos chamar $S(\omega)$ de $\omega+1$, $S(S(\omega))$ de $\omega+2$, e, de forma geral, $S^n(\omega)$ de $\omega+n$.
\[
0, 1, 2, \dots, \omega, \omega+1, \omega+2, \dots
\]

Tal processo de fato produz boas ordens: pode-se verificar diretamente que $\omega+1, \omega+2, \omega+3$ são conjuntos bem ordenados por $\in$.
Assim, uma primeira ideia é simplesmente definir um ordinal como um conjunto bem ordenado por $\in$.
No entanto, a unicidade de representantes de boas ordens não é garantida por essa definição.
Considere os dois conjuntos abaixo:
\[
    A=5=\{0, 1, 2, 3, 4\}
\]
\[
    B=\{0, 1, 2, 3, 6\}
\]
Ambos são conjuntos bem ordenados por $\in$ (pois são finitos), e são isomorfos, mas são distintos.
O ``truque'' utilizado para construir $B$ foi remover de $7$ um elemento que ``deveria estar lá''.
Em termos mais precisos, existe um $a \in 6$ tal que $a \notin B$ (no caso, $a=5$).
Em fórmula de primeira ordem:

\[
    \exists b\in B \exists a\in b (a\notin B).
\]

Optaremos por proibir esse ``truque''.
A negação de tal fórmula é a seguinte:
\[
    \forall b\in B \forall a\in b (a\in B).
\]
Ou seja:
\[
    \forall b\in B (b\subseteq B).
\]

Conforme veremos, com a presença do Axioma da Substituição, que será introduzido na próxima seção, conjuntos bem ordenados por $\in$ que satisfazem essa condição são exatamente os representantes de todas as boas ordens.

\begin{definition}
    Um conjunto $A$ é dito \emph{transitivo} se, para todo $a\in A$, $a\subseteq A$.
\end{definition}

Agora podemos definir um ordinal.

\begin{definition}
    Um \emph{ordinal} é um conjunto $A$ tal que:
    \begin{enumerate}[label=(\alph*)]
        \item $A$ é bem-ordenado por $\in$ (ou seja, $(A, \in_A)$ é uma boa ordem);
        \item $A$ é transitivo.
    \end{enumerate}
\end{definition}
No contexto de ordinais, escrevemos sempre que for conveniente, $a<b$ para $a\in b$.

\begin{example}
    $\omega$ é um ordinal: já vimos que $\omega$ é um conjunto bem-ordenado por $\in$, e, para todo $a\in \omega$, $a$ é um número natural, e, portanto, $a\subseteq \omega$.
\end{example}

Além disso, todos os números naturais são ordinais pelo lema abaixo.

\begin{lemma}
    Seja $\alpha$ um ordinal.
    Então todo elemento de $\alpha$ é um ordinal.
\end{lemma}
\begin{proof}
Seja $\beta \in \alpha$.
Temos que $\beta\subseteq \alpha$, logo, $\beta$ é um conjunto bem-ordenado por $\in$.
Resta ver que $\beta$ é transitivo.
Seja $\gamma \in \beta$.
Vejamos que $\gamma\subseteq \beta$.
Para isso, tome $\delta \in \gamma$.
Temos que $\delta \in \gamma \in \beta \in \alpha$.
Como $\alpha$ é transitivo, $\delta \in \gamma \in \alpha$.
Novamente, como $\alpha$ é transitivo, $\delta \in \alpha$.
Assim, $\delta<\gamma<\beta$ em $\alpha$, e, portanto, $\delta <\beta$, ou seja, $\delta \in \beta$.
\end{proof}

A relação $\in$ é, em verdade, uma boa ordem na classe dos ordinais.

Precisaremos de alguns lemas.

\begin{lemma}
    Seja $A$ um conjunto não vazio de ordinais.
    Então $\bigcap A$ é um ordinal.
\end{lemma}
\begin{proof}
    Seja $\alpha\subseteq A$ qualquer.
    Então $\bigcap A\subseteq \alpha$, logo, $\bigcap A$ é um conjunto bem-ordenado por $\in$.

    Resta ver que $\bigcap A$ é transitivo.
    Seja $\beta \in \bigcap A$ e $\gamma \in \beta$.
    Para todo $\xi \in A$, $\beta \in \xi$, logo, $\gamma \in \xi$ (pois $\xi$ é transitivo).
    Assim, $\gamma \in \bigcap A$.
\end{proof}

Dados $\alpha, \beta$ ordinais, pela nossa convenção de escrita, temos que $\alpha<\beta$ se, e somente se $\alpha\in \beta$.

O seguinte lema nos dá uma equivalência para $\leq$.
\begin{lemma}
    Sejam $\alpha, \beta$ ordinais.
    Então $\alpha\subseteq \beta$ se, e somente se $(\alpha=\beta \land \alpha\in \beta)$.
\end{lemma}
\begin{proof}
    Se $\alpha=\beta$, é imediato que $\alpha\subseteq \beta$.
    Se $\alpha\in \beta$, então $\alpha\subseteq \beta$, pois $\beta$ é transitivo.

    Reciprocamente, suponha que $\alpha\subseteq \beta$.
    Suponha ainda que $\alpha\neq \beta$.
    Veremos que $\alpha\in \beta$.

    Seja $X=\beta\setminus \alpha$ e $\delta=\min X$.
    Veremos que $\delta=\alpha$, o que completa a prova uma vez que $\delta\in \beta$.

    Ora, temos que se $\xi<\delta$, então $\xi\notin X$.
    Como $\xi\in \beta$, então $\xi\in \alpha$.
    Isso nos mostra que $\delta\subseteq \alpha$.

    Suponha por absurdo que existe $\gamma \in \alpha\setminus \delta$.
    Então $\gamma, \delta \in \beta$ (pois $\beta$ é transitivo), logo $\gamma$ e $\delta$ são comparáveis.
    
    Como $\gamma\notin \delta$, temos que $\delta\in \gamma$ ou $\delta=\gamma$.
    $\delta=\gamma$ é impossível, pois $\delta\notin \alpha$ e $\gamma\in \alpha$.
    Por outro lado, $\delta \in \gamma$ é impossível, pois teríamos $\delta\in \gamma \in \alpha$, e, portanto, $\delta\in \alpha$, o que é absurdo.

    Assim, $\alpha\subseteq \delta$, e, portanto, $\alpha=\delta$.
\end{proof}
\begin{proposition}
    Sejam $\alpha, \beta, \gamma$ ordinais.
    Então:

    \begin{enumerate}[label=(\alph*)]
        \item $\alpha\notin \alpha$;
        \item se $\alpha\in \beta$ e $\beta\in \gamma$, então $\alpha\in \gamma$;
        \item $\alpha\in \beta$ ou $\alpha=\beta$ ou $\beta\in \alpha$.
        \item Todo conjunto não vazio de ordinais tem um elemento mínimo.
        Ainda mais, se $X$ é um conjunto não vazio de ordinais, $\bigcap X=\min X$.
    \end{enumerate}        
\end{proposition}
\begin{proof}
    (a) Se $\alpha\in \alpha$, então $\exists x \in \alpha(x \in x)$, o que é impossível, pois $\in$ é irreflexivo em $\alpha$.

    (b) Suponha que $\alpha\in \beta$ e $\beta\in \gamma$.
    Como $\gamma$ é transitivo, $\beta\subseteq \gamma$, logo, $\alpha\in \gamma$.

    (c) Sejam dados $\alpha, \beta$ ordinais.
    Temos que $\gamma=\alpha\cap \beta$ é um ordinal, e $\gamma\subseteq \alpha$ e $\gamma\subseteq \beta$.

    Se $\gamma=\alpha$, temos que $\alpha\subseteq \beta$, logo, $\alpha=\beta$ ou $\alpha\in \beta$.

    Se $\gamma=\beta$, temos que $\beta\subseteq \alpha$, logo, $\alpha=\beta$ ou $\beta\in \alpha$.

    Se $\gamma\subsetneq \alpha$ e $\gamma\subsetneq \beta$, então $\gamma\in \alpha$ e $\gamma\in \beta$, logo, $\gamma\in \alpha\cap \beta=\gamma$, o que é absurdo.

    (d) Seja $X$ um conjunto não vazio de ordinais e $\alpha=\bigcap X$.
    Para todo $\beta\in X$, $\alpha\subseteq \beta$, logo, $\alpha\leq \beta$ para todo $\beta\in X$.
    Resta ver que $\alpha \in X$.
    Ora, para todo $\delta \in X$, temos que $\alpha\subseteq \delta$.
    Se $\alpha\notin X$, então $\alpha\in \delta$ para todo $\delta \in X$, logo, $\alpha\in \bigcap X=\alpha$, o que é absurdo.
\end{proof}

\begin{corollary}
Todo conjunto de ordinais é bem-ordenado por $\in$.
Assim, um conjunto de ordinais transitivo é um ordinal.
\end{corollary}
\begin{corollary}
Não existe $X$ tal que $\forall \alpha(\alpha \text{ é ordinal} \rightarrow \alpha\in X)$.
\end{corollary}
\begin{proof}
    Se existisse, seja $O=\{\alpha \in X: \alpha \text{ é ordinal}\}$.
    Então $O$ é o conjunto de todos os ordinais, e, portanto, é bem-ordenado por $\in$.
    Como todo elemento de um ordinal é um ordinal, $O$ é transitivo, logo, $O$ é um ordinal.
    Assim, $O\in O$, o que é absurdo.
\end{proof}
\begin{lemma}
    Para todo ordinal $\alpha$, então $S(\alpha)=\alpha\cup\{\alpha\}$ é um ordinal.
\end{lemma}

\begin{proposition}
    Seja $\alpha$ um ordinal.
    Então $S(\alpha)$ é um ordinal.
    Mais especificamente, $S(\alpha)$ é o menor ordinal $\beta$ tal que $\alpha<\beta$.
\end{proposition}
\begin{proof}
    $S(\alpha)$ é um conjunto de ordinais, então basta ver que ele é transitivo.
    Se $\beta \in S(\alpha)$, então $\beta\in \alpha$ ou $\beta=\alpha$.
    Em qualquer caso, $\beta\subseteq \alpha\subseteq S(\alpha)$.
    Assim, $S(\alpha)$ é um ordinal.

    Obviamente, $\alpha \in S(\alpha)$.
    Finalmente, se $\gamma<S(\alpha)$ é um ordinal, então $\gamma\leq \alpha$, assim, não existe $\gamma$ tal que $\alpha<\gamma<S(\alpha)$.
\end{proof}

A classe de todos os ordinais é denotada por $\ON$.
Formalmente, $\ON$ não é uma relação unária, e escrevemos $\alpha \in \ON$ como uma outra forma de escrever $\ON(\alpha)$, que se traduz como ``$\alpha$ é um ordinal''.


\begin{lemma}
    Seja $X$ um conjunto de ordinais.
    Então $\bigcup X$ é um ordinal.
    Além disso, $\bigcup X$ é o supremo de $X$ em $\ON$, ou seja, o menor ordinal $\alpha$ tal que $\beta\leq \alpha$ para todo $\beta\in X$.
\end{lemma}
\begin{proof}
    Como $X$ é um conjunto de ordinais, e todo elemento de um ordinal é um ordinal, segue que $\bigcup X$ é um conjunto de ordinais.
    Temos que $\bigcup X$ é transitivo, pois se $\alpha \in \bigcup X$, então $\alpha\in \beta$ para algum $\beta\in X$, e, portanto, $\alpha\subseteq \beta\subseteq \bigcup X$.

    Assim, $\bigcup X$ é um ordinal e, para todo $\alpha \in X$, $\alpha\subseteq \bigcup X$, ou seja, $\alpha\leq \bigcup X$.

    Finalmente, se $\beta<\bigcup X$, temos que existe $\alpha \in X$ tal que $\beta\in \alpha$, logo, $\bigcup X$ é a menor cota superior de $X$.
\end{proof}

Nem todo ordinal é o sucessor de outro ordinal.
Por exemplo, $\omega$ não é o sucessor de nenhum ordinal.

\begin{definition}
    Um ordinal $\beta$ é dito \emph{sucessor} se é da forma $\beta=S(\alpha)$ para algum ordinal $\alpha$.

    Um \emph{ordinal limite} é um ordinal que não é $0$ nem sucessor.
\end{definition}

\begin{proposition}
    $\omega$ é um ordinal limite.
\end{proposition}
\begin{proof}
Se existisse $\alpha$ tal que $S(\alpha)=\omega$, teríamos que $\alpha\in \omega$.
Porém, isso implica que $\omega=S(\alpha)=\omega$, um absurdo.
Por outro lado, $\omega\neq 0$, pois $0\in \omega$ e $0\notin 0$.
\end{proof}

\subsection*{Exercícios}
\begin{exercise}
    Prove que a interseção de uma coleção não vazia de conjuntos transitivos é um conjunto transitivo.
\end{exercise}
\begin{exercise}
    Prove que a união de uma coleção de conjuntos transitivos é um conjunto transitivo.
\end{exercise}
\begin{exercise}
    Prove que $\omega$ é o menor ordinal limite, ou seja, que $\omega$ é um ordinal limite e, para todo ordinal limite $\alpha$, $\omega\leq \alpha$.
\end{exercise}